% Drift-Aware Adaptive Page Replacement -- LaTeX transcription of the revised manuscript.
% Baseline was a FAITHFUL transcription of the now-removed historical PDF Drift_Aware_Page_Replacement_Revised.pdf; this source is authoritative.
% Fixes are then applied one by one; see CHANGELOG.md.
\documentclass[10pt,twocolumn,a4paper]{IEEEtran}
\usepackage[T1]{fontenc}
\usepackage[utf8]{inputenc}
\usepackage{amsmath,amssymb}
\usepackage{graphicx}
\usepackage{booktabs}
\usepackage{array}
\usepackage{makecell}
\usepackage{xcolor}
\usepackage[hidelinks]{hyperref}
\usepackage{balance}
\graphicspath{{figs/}}
\setlength{\tabcolsep}{4pt}
\renewcommand{\arraystretch}{1.0}

\begin{document}

\title{When Does Drift Gating Pay Off? A Reproducible Measurement Study of Drift-Aware Page Replacement}

\author{Muhammad~Qasim,~Hunain~Tahir,~Hussain~Ahmad,~and~Mohsin~Ghaffar~Ghouri%
\thanks{M. Qasim, H. Tahir, H. Ahmad, and M. G. Ghouri are with the University
of Management and Technology (UMT), Lahore, Pakistan. M. Qasim is also with the
Anoneurx Research Lab, Anoneurx. Corresponding author: M. Qasim (e-mail:
itskashie@anoneurx.com).}}

\maketitle

\begin{abstract}
When drift-triggered gating pays for adaptive page replacement is unclear. We
report a reproducible measurement study of a CUSUM trigger with sampled
shadow evaluation and a cost-margin switch; an active KS confirmation stage had
small ablation effect and is not a validated gate. On 80 non-stable synthetic
traces, gating saved 76\% of shadow evaluations versus continuous evaluation
(6.47 vs.\ 27.39/trace) but had higher miss ratio (0.453 vs.\ 0.423) and was
not significantly different from S3-FIFO (0.468). On six public FIU block
traces, the frozen controller never switched because of an identified
capacity-sampling fault; a post hoc repair produced identical public,
synthetic, and holdout results. Post hoc follow-ups found a sharp
shadow-rankability boundary ($C_s/D_w<1$) but no meaningful FIU repair gain;
CUSUM dominated periodic/random at matched budgets, while Oracle/ContShadow
exposed an accuracy ceiling; selected Page-Hinkley and simplified ADWIN
detectors did not close the oracle gap. On available real traces and added
SIEVE/W-TinyLFU baselines, DA showed no accuracy advantage. A user-space
minor-fault calibration places the DA--ContShadow crossover at 2.2\,$\mu$s
(95\% CI 1.44--4.06\,$\mu$s), but a total-time envelope that includes CLOCK and
S3-FIFO shows DA is never the cheapest policy on the measured means: CLOCK is
cheapest below 0.216\,$\mu$s, S3-FIFO between 0.216 and 2.305\,$\mu$s, and
continuous evaluation above 2.305\,$\mu$s, with DA at least 3.27\,ns per
reference above the best alternative at every miss cost. Gating reduces shadow
work, but no accuracy advantage was observed in these measurements, and DA is
never best in these measurements.
\end{abstract}

\section{Introduction}
\label{sec:intro}

Physical memory is finite. When a process references a page that is not
resident and no free frame is available, the memory manager must evict a
resident page. The rule that makes this choice, the page-replacement policy,
affects page-fault frequency, storage traffic, tail latency and, ultimately,
application throughput. Recency-based policies such as least recently used
(LRU) and its CLOCK approximation remain the practical default in operating
systems, while the research literature offers frequency-aware, scan-resistant
and self-tuning alternatives, including the adaptive replacement cache (ARC)
[1], CLOCK with adaptive replacement (CAR) [2], the low inter-reference recency
set policy (LIRS) [3], 2Q [4] and S3-FIFO [5].

No single policy is best for every access pattern. A sequential scan can
flush a cache tuned for repeated reuse, a loop slightly larger than memory can
make LRU miss on every reference, and a policy that protects a previously hot
working set can retain pages that are useless after a phase change. Real
workloads move between such regimes: a database alternates between index
lookups and table scans, a scientific code between computation kernels, and a
multitasking system between applications. The question this paper studies is
therefore not which policy is best, but whether a system can recognize that
the reference behaviour has changed and respond at an acceptable cost.

The stakes of this choice have grown. Containerized and multi-tenant servers
interleave unrelated workloads on shared memory, GPU and accelerator memory is
oversubscribed by multitasking and large-model serving, and tiered memory
systems make a wrong eviction more expensive because the next level is slow or
recomputation is costly [26]--[28]. Recent studies show that adaptive and
learned policies can reduce misses substantially [15], [18], [19], but also
that simple policies such as LRU remain surprisingly competitive in some
settings [28] and that policy rankings depend strongly on the workload [29]. A
principled account of when adaptation pays for itself is therefore needed, not
only another adaptive algorithm.

We use workload drift to mean a statistically meaningful change over time in
the distribution or temporal structure of page references. Two research lines
address related problems. Adaptive replacement algorithms infer from feedback
which behaviour to favour: ARC shifts capacity between recency and frequency
lists, LeCaR and CACHEUS reweight expert policies by regret, and recent learned
or learning-augmented designs reduce the cost of learning on the eviction path
[11], [12], [15], [19]. Concept-drift research, in contrast, provides formal
change detectors with controllable error rates, such as adaptive windowing and
sequential tests [6], [7], [20]--[25]. The two lines have overlapping goals
but different constraints. A drift detector built for a machine-learning
pipeline may be too expensive, or rest on assumptions (labels, prediction
errors) that do not exist in a page-reference stream. A replacement algorithm
that adapts continuously may, conversely, switch behaviour without any control
over the rate of unnecessary switches.

A closely related precedent exists outside operating-system paging. Adaptive
cache-mode selection for queries over raw data [34] keeps a shadow cache of the
alternative mode and switches when a significance test favours it. That work
shows that statistically guided cache adaptation is not new; what remains open
is whether drift-gated selection among page-replacement policies is effective
once the costs of monitoring, shadow evaluation and policy-state migration are
counted, and how it compares with strong existing replacement algorithms.

Four constraints make drift-aware replacement hard. (a)~Rate. References
arrive at memory speed, so any per-reference statistical work is prohibitive,
and monitoring must be sampled or summarized by epoch. (b)~Missing labels and
policy-dependent signals. There is no ground truth for ``the workload
changed'', and the most obvious signal, the miss ratio, depends on the active
policy and capacity, so a switch can itself create an apparent change.
(c)~State. Policies keep different metadata, and a switch must rebuild it from
the resident pages; the quality of that migration, and which pages it retains,
can decide whether a switch helps or harms. (d)~Error control. False alarms
waste work and cause oscillation, while missed changes leave a mismatched
policy in place. A design that ignores any of the four can look good on paper
and fail in operation.

This paper makes three contributions. (1)~A controlled measurement of drift
gating against continuous shadow evaluation and strong static baselines, on
freshly generated phase traces, on six public block traces, and on
instrumented real programs with whole programs held out. We report the null
and negative outcomes: the gate is not significantly different from S3-FIFO,
continuous evaluation is more accurate, and on the public traces the gated
policy never switches. (2)~A cost break-even model for the shadow-work saving,
calibrated by a measured user-space minor-fault latency, together with a
total-time envelope over CLOCK, S3-FIFO, ContShadow and DA showing that DA is
never best in these measurements (Section~\ref{sec:breakeven}). (3)~A fully reproducible
negative-result protocol: released simulator, seeds, configurations, per-trace
CSVs, selection sweep, ablations, sensitivity checks, and scripts that
regenerate every table and figure. We do not claim a new replacement algorithm
or a validated statistical gate; the effective mechanism is a CUSUM trigger
with sampled shadow evaluation.

Our research questions are: RQ1 How do recent adaptive replacement policies
detect and respond to changing access behaviour? RQ2 Which drift detectors are
candidates for low-overhead use in memory management? RQ3 What do emerging
GPU, multi-tier and benchmark studies imply for replacement under changing
workloads? RQ4 Can a drift-gated controller with sampled shadow evaluation
select among replacement policies more effectively than fixed and
heuristic-adaptive baselines under controlled phase changes? RQ1--RQ3 are
answered by the review;
RQ4 is answered empirically.

We adopt a deliberately cautious evaluation stance. The protocol fixes the
design and parameters on training traces before test traces are touched,
evaluates all methods on identical traces, uses paired non-parametric
statistics with multiple-comparison correction, includes an offline optimum so
that the remaining gap is visible, and reports negative and null findings. We
also compare against an Acme-style variant that applies the same shadow
evaluation without a drift gate [34], because that comparison isolates what
drift detection itself contributes.

The answer to RQ4 is mixed, and we state it here so that the scope of the
paper is clear. On 80 fresh non-stable generated traces the gated controller
reduced the mean miss ratio by 20.5\% relative to LRU, but was not
significantly different from S3-FIFO (Holm-adjusted $p = 0.09$) and was
significantly worse than continuous shadow evaluation by 0.030 miss ratio. It
cut shadow evaluations by 76\% relative to continuous evaluation, but in this
study not at equal miss ratio. On six public block traces the frozen controller
never switched because of an identified capacity-sampling fault and was
Pareto-dominated by S3-FIFO, a failed external validation; a post hoc repaired
variant selected the same configuration and produced identical public, synthetic, and holdout results. On 15 instrumented programs with whole programs
held out it tied S3-FIFO at higher cost, and no comparison survived Holm
correction. A cost model puts the DA--ContShadow crossover at about $2.2\,\mu$s per miss,
close to the measured minor-fault latency, but including CLOCK and S3-FIFO in
the total-time envelope shows DA is never best in these measurements: CLOCK,
S3-FIFO and ContShadow each occupy part of the lower envelope. We therefore
report where gating saves shadow work, not superiority over existing methods
(Section~V).

The evidence has clear limits. The main experiments use generated traces --
phase traces with known change points and composites of five access models --
not traces collected from running systems. The public block traces are
storage/cache traces rather than virtual-memory paging traces and serve as an
external check only. The Appendix reports a separate, exploratory evaluation
on page traces from 15 instrumented programs with whole programs held out: DA
tied S3-FIFO there, and no comparison survives Holm correction. A faithful
kernel multi-generation LRU (MGLRU) baseline is not possible in page-level
trace replay because the kernel aging and reclaim-pressure signals are absent;
a page-level MGLRU-style abstraction is reported but not claimed as kernel
MGLRU. No production paging traces and no in-kernel page-fault or reclaim
measurements were collected; we do report a user-space minor-fault latency
calibration for the cost model. The simulator measures algorithmic behaviour
and relative cost, not kernel performance. Section~VI and Section~VII discuss
these limitations in detail.

The remainder of the paper is organized as follows. Section~II gives
conceptual background. Section~III states the review scope and synthesizes the
reviewed studies, and Section~IV identifies the research gaps. Section~V
presents the experimental methodology and results. Section~VI discusses
threats to validity, Section~VII the limitations and scope of the claims, and
Section~VIII concludes.

\section{Conceptual Background}
\label{sec:background}

\subsection{Page Replacement and Workload Phases}

A reference trace is an ordered sequence of page identifiers. For a capacity of
$C$ frames, a replacement policy chooses the victim whenever a reference misses
and all frames are occupied. LRU evicts the least recently used page; LFU the
least frequently used; FIFO the oldest; CLOCK approximates LRU with a reference
bit and a circular scan; MRU evicts the most recently used page, which is
favourable for cyclic scans larger than memory. Hybrid designs use several
queues to separate one-time from repeated references [1]--[5]. The miss ratio
is the fraction of references that miss; it is a necessary but incomplete
quality measure, since metadata cost, decision cost and adaptation speed also
matter. Results on storage-block traces are cache-replacement results and are
not direct evidence about virtual-memory paging.

Workload phases differ in working-set size, reuse-distance distribution,
sequentiality and frequency skew. Drift may be abrupt, gradual or recurring
[6]. Phase changes also interact with policy state: a policy that retains
pages of the previous phase (for example MRU, which pins the oldest resident
pages) can perform far worse than its steady-state behaviour suggests until
that state is flushed. This interaction is central to the design and results
of this paper.

\subsection{Workload Drift and Statistical Detection}

Concept-drift research distinguishes changes in the input distribution, in the
input--output relationship, or both [6]. A page-reference stream has no
labels, so a controller must monitor observable summaries such as the miss
ratio or reuse-time distribution. Adaptive windowing (ADWIN) tests whether
sub-windows have different means [7]; Page--Hinkley and CUSUM procedures
accumulate deviations from a reference level and signal when the sum exceeds a
threshold [30]; two-sample tests such as the Kolmogorov--Smirnov (KS) test
compare distributions. Each has assumptions about sample size, independence
and repeated testing that must be addressed in an implementation. Because
memory operations occur at very high rates, a detector for this setting must
avoid per-reference heavy computation and must balance detection delay, false
alarms, memory use and runtime cost.

\subsection{Earlier Adaptive Replacement and Cache Modelling}

SEQ [8] detects long fault sequences and handles them separately, and EELRU
[9] adjusts its eviction point from recency-distribution evidence. ARC [1]
balances recency and frequency using ghost-list hits, CAR [2] transfers the
idea to CLOCK, LIRS [3] uses inter-reference recency, and 2Q [4] separates
newly admitted from re-referenced pages. S3-FIFO [5] shows that FIFO queues
with a small admission queue can match or beat more complex designs. ACME
[10], LeCaR [11] and CACHEUS [12] mix expert policies using regret or learned
feedback; Acme [34] is the closest precedent for statistically guided
switching, but it selects a query-caching mode rather than a page-replacement
policy. For evaluating alternatives cheaply, miniature simulation [13]
emulates a larger cache with a smaller one, and SHARDS [14] estimates
miss-ratio curves from spatially hashed reuse-distance samples; both make
on-demand shadow evaluation feasible, although sampling can distort temporal
structure and policy rankings. On the change-detection side, ADWIN [7] tests
whether sub-windows of an unlabelled stream have different means; it is a
candidate low-cost trigger but has not been evaluated at page-reference rates.
Set dueling in DIP/DRRIP [37] and the multi-generation LRU reclaim path [38]
are positioned in Section~III. These works establish that adaptation,
sampling and change detection each have precedent; the controller studied
here combines them and measures the consequences, and makes no priority claim
for the combination itself.

\section{Review Scope and Approach}
\label{sec:review}

This is a focused narrative review, not a systematic review or meta-analysis.
The study set was selected by the authors to cover adaptive replacement, drift
detection and emerging memory or cache settings, and comprises 15 peer-reviewed
publications and preprints from 2025--2026 plus foundational work needed to
explain the lineage. One study [20] first appeared as a preprint in December
2024 and was published at AAAI 2025. No preregistered search protocol,
duplicate screening or quality scoring was applied, so the review does not
claim systematic completeness, and several included works are preprints.
Results reported in the studies are not pooled or normalized because
workloads, hardware and metrics differ.

\subsection{Adaptive and Learning-Based Replacement}

3L-Cache [15] reduces the cost of learned eviction by moving learning and
inference off the critical path, evaluated on 4,855 CDN, key-value and block
traces, but its overhead remains above LRU. LPR [16] uses re-access feedback to
switch among policies at runtime, for example from LRU to MRU, with a
heuristic rather than statistical switching signal. EACLOCK [17] adapts a
CLOCK-style policy from hit-rate feedback for skewed and shifting workloads.
Clock2Q+ [18] targets correlated metadata references in VMware vSAN and reports
up to 28.5\% lower miss ratio than S3-FIFO on metadata traces. S4-FIFO [19]
learns configuration parameters for a FIFO-family heuristic rather than every
eviction, and, like 3L-Cache, separates a fast data path from slower control
logic.

\subsection{Drift Detection}

PUDD [20] detects drift from prediction uncertainty but needs a predictive
model. A label-free statistical-process-control detector [21] stresses power
and false-positive control but is evaluated by simulation of model-monitoring
scenarios. Dynamic threshold determination [22] adapts alarm thresholds after
each alarm, at the cost of maintaining parallel candidate models temporarily.
A transformer--autoencoder hybrid [23] and an HTM-plus-SPRT framework [24]
offer richer temporal modelling but impose training or univariate-data
assumptions. ADWIN-U [25] adapts adaptive windowing to unsupervised settings.
None is evaluated on page-reference streams or at memory-management rates.

\subsection{GPU, Multi-Tier and Benchmark Studies}

MSched [26] predicts the next task's working set from kernel launch arguments
to migrate GPU memory proactively, which relies on predictability that CPU page
streams lack. gpu\_ext [27] makes GPU-driver policies replaceable eBPF
programs, an architectural mechanism for safe policy updates. The LLM
prefix-cache study [28] finds LRU surprisingly competitive, cautioning against
assuming that sophisticated policies always win. A stale-tolerant cache
benchmark [29] finds that the ranking is workload-dependent across four
synthetic patterns, with ARC highest on average. Finally, Acme [34] is the closest precedent for statistical policy
switching: it simulates the alternative caching mode through shadow entries and
switches when a one-tailed $t$-test at the 0.05 level favours it, reporting
execution time within 16\% of the optimal mode in the worst case and 4\% on
average on TPC-H workloads.

Two further lines of prior art are directly relevant but were absent from the
original 15-study review. Set dueling, introduced for dynamic insertion
policies (DIP/DRRIP) by Qureshi \emph{et al.} [37], dedicates sampler sets to
two candidate insertion policies and lets the one with fewer sampler misses
control the whole cache at very low storage cost; it is the standard
low-overhead way to choose between replacement behaviours. In deployed
operating systems, the multi-generation LRU (MGLRU) reclaim path [38] groups
pages by access recency into generations and ages them incrementally, replacing
the single active/inactive list. Both are sampled or hierarchical ways of
adapting recency treatment, and both are therefore natural comparators for a
drift-gated controller. These descriptions are direct observations from the
cited sources. That the mechanisms would transfer to drift-gated page
replacement is \emph{inferred}: none of the cited evaluations measures shadow
evaluation cost, policy-state migration and post-switch warm-up together on
page-reference streams. Section~VII and the Appendix add principle-level
versions of them; the main tables predate this addition and do not include
them.

\section{Previous Research Gaps}
\label{sec:gaps}

The 15 reviewed studies and the closest prior art are summarized in Table~S.I
(supplementary material) by research area, contribution and limitation relative
to low-overhead, drift-aware page replacement. The limitations are framed
against that goal; they are not claims that the original studies fail at their
own objectives, and some are inferred from each study's scope. Reported
results are not recomputed or compared across studies.



The reviewed work establishes that adaptive replacement is well studied
[1]--[5], [8]--[12]; that learning-based policies can respond online to
access patterns [15]--[19]; that statistical drift detection offers explicit
change evidence [6], [7], [20]--[25]; and that cache-mode selection has used
shadow evaluation and significance testing [34]. Four gaps follow.

Gap~1: detectors on reference streams. Generic drift detectors are rarely
evaluated on page-reference feature streams at memory-management sampling
rates and overhead budgets. Gap~2: statistical gating versus heuristic
adaptation. Adaptive page-replacement studies such as LPR [16] already respond
to changing patterns, but no reviewed study compares explicit statistical
change confirmation against such heuristics in a controlled experiment.
Gap~3: transition costs. The costs of detector execution, shadow simulation,
policy-state preparation and post-switch warm-up are not consistently reported
together with miss ratio and detection delay. Gap~4: comparability. Studies use
different workloads, metrics and system contexts, so their results cannot be
compared as though from a shared benchmark. Gaps~2--4 are direct observations
about the reviewed set. Gap~1 is inferred from the stated scopes of the cited
detectors; it is not a claim that no such evaluation exists outside this set.

Table~S.II (supplementary material) compares the reviewed approaches with
the controller as actually implemented in this study.



These gaps are empirical. The question is whether statistically confirmed
changes in compact page-reference summaries, combined with on-demand candidate
evaluation and explicit switching-cost safeguards, give a favourable trade-off
against fixed policies and existing adaptive methods. Each component of the
proposed controller has precedent (ADWIN [7], miniature simulation [13],
SHARDS [14], statistical cache-mode selection [34]); the claim examined here
concerns their integration for page-replacement policy selection and the
measured consequences. The review cannot establish that no related
implementation exists outside the selected literature, and the proposal is not
presented as the first statistical cache-adaptation system.

\section{Experimental Methodology and Results}
\label{sec:eval}

\emph{Reproducibility.} Every experimental number in this paper is produced
by released code in the accompanying package: the synthetic trace generator,
all baselines, the controller, the seeds and the raw per-trace output, the
parameter-selection sweep (\texttt{sweep\_selection.py}), the primary-study
driver (\texttt{run\_main\_study.sh}), the table builder
(\texttt{analyze\_main\_study.py}), the appendix interval script
(\texttt{appendix\_intervals.py}), the cost model (\texttt{cost\_model.py}),
the public-trace converter (\texttt{convert\_public\_traces.py}) and the
fault-latency microbenchmark (\texttt{prototype\_fault\_latency.cpp}) are
provided. The selection sweep is run with the released simulator on training
traces (seed 100), and the selected configuration is evaluated once on fresh
traces (seed 5001).

An earlier report of this study used a simulator, seeds and raw output that
were not retained; its tables cannot be regenerated, none of its numbers is
used here, and it is not discussed further.
\begin{figure}[!t]
\centering
\includegraphics[width=0.85\columnwidth]{fig2_workflow.png}
\caption{Experimental workflow. Blue boxes are training stages, orange boxes
are held-out test rounds.}
\label{fig:workflow}
\end{figure}

\subsection{Simulator, Workloads and Baselines}

We implemented a single-threaded trace-driven simulator in C++ with a capacity
of $C = 1{,}000$ frames, excluding the first 20,000 references (cold fill)
from reported miss ratios. All policies see identical traces. Policies were
sanity-checked against known behaviour: LRU misses on every reference of a
loop larger than memory, and MRU matches the offline optimum on that loop. OPT
minimises misses over the whole trace including the cold fill, so it is a
reference point rather than a strict lower bound on the post-warm-up miss
ratio. Workloads (Table~III) are generated, not
collected: five access primitives with disjoint page ranges, combined into
stable, abrupt, gradual and recurring scenarios with known change points, plus
composite traces from five kernel-level access models (column-wise matrix
traversal, hash join, B-tree lookup, external-sort merge and skewed key-value
lookup with scan bursts). Traces are 110 per round: 30 stable, 20 abrupt,
20 gradual, 20 recurring and 20 composite.

\begin{table}[!t]
\caption{Workload Definitions}
\label{tab:workloads}
\footnotesize
\begin{tabular}{@{}p{0.24\columnwidth}p{0.70\columnwidth}@{}}
\toprule
Component & Definition \\
\midrule
LOOP & Cyclic scan over 1,200 pages (larger than $C$). \\
ZIPF & Zipf(0.9) over 8,000 pages, randomly permuted page ids. \\
HOTSCAN & 70\% uniform over 600 hot pages, 30\% sequential scan over 20,000 pages. \\
WINDOW & Uniform over a 700-page window sliding one page per 40 references. \\
UNIFORM & Uniform random over 4,000 pages. \\
Stable & One primitive, 300,000 refs (6 traces each). \\
Abrupt & Six phases of 100,000 refs; no immediate repetition of a type. \\
Gradual & As abrupt, with linear mixture ramps of 40,000 refs at boundaries. \\
Recurring & Two primitives alternating, $6 \times 80{,}000$ refs. \\
Composite & Five kernel models in random order, $5 \times 100{,}000$ refs. \\
\bottomrule
\end{tabular}
\end{table}

Table~VI reports eleven non-proposed baselines (eight fixed and three
adaptive), while Table~VII pairs DA against ten of them because MRU is omitted
from the paired tests; the counts of eleven and ten refer to different tables,
not to different sets of results.

The two count-aging periods reported in Tables~IV and~V apply to different
counters, and both are intended. Table~IV describes the LFU-A policy itself:
its frequency counts are halved every $4C$ references (4,000 at
$C = 1{,}000$), and the same rule governs the LFU-A candidate inside shadow
evaluation, scaled to the capacity in use (every $4C_s = 800$ references at
$C_s = 200$). Table~V's ``Count aging'' row instead describes the controller's
own fast-path aged count (Section~V-C), which is halved every 50 epochs, i.e.\
every 100,000 references; that counter is not part of LFU-A and is used,
together with the last-access time, only to place pages when a switch migrates
state.

Re-implementations of published
algorithms follow the cited descriptions but are not the authors' code; the
LPR-style baseline re-implements the principle of [16] (re-access ``regret''
feedback with probationary LRU/MRU switching), not its code, and may differ
from the original in important ways. CAR, LIRS, CLOCK-Pro [36], CACHEUS and a
3L-style policy are included in the released simulator and raw CSVs but are not
reported in the main tables; S4-FIFO, EACLOCK and Clock2Q+ were not implemented.
The ContShadow variant is our Acme-style control: it runs the same shadow
evaluation and switching rule every ten epochs without a drift gate.

\begin{table}[!t]
\caption{Methods Compared}
\label{tab:methods}
\footnotesize
\begin{tabular}{@{}p{0.16\columnwidth}p{0.78\columnwidth}@{}}
\toprule
Group & Methods \\
\midrule
Fixed & LRU, FIFO, CLOCK, LFU-A (counts halved every $4C$ references; the rule scales with the capacity in use), MRU, 2Q ($K_{\mathrm{in}} = C/4$, $K_{\mathrm{out}} = C/2$), ARC [1], S3-FIFO [5] (10\% small queue, ghost = main size). \\
Adaptive & LeCaR [11] (re-implementation); LPR-style (principle of [16]); ContShadow (shadow evaluation every 10 epochs, no drift gate). \\
Proposed & DA: drift-gated controller (Section~V-C). \\
Ablations & DA-noPurge (A0 migration), DA-noS2 (no KS confirmation), DA-noShadow (random pick), DA-noCost (no cost margin), DA-noDwell. \\
Reference & OPT: offline Belady optimum [32]. \\
\bottomrule
\end{tabular}
\end{table}

The protocol has two stages. (i)~\emph{Training} (seed 100, 15 stable and 32
non-stable traces). A sweep over candidate-set size $\{4,5,6\}$, cost margin
$K \in \{0, 0.05, 0.1, 0.2\}$ and evaluation delay $\{2,5,10,20\}$ selected
five candidates (adding S3-FIFO), $K = 0$, and evaluation delay 10 by the lowest
training miss ratio with zero stable-trace false confirmations. (ii)~\emph{Test}
(fresh seed 5001, 30 stable and 80 non-stable traces), evaluated once with the
frozen configuration. No parameter was changed after the test run. The sweep is
re-executed by \texttt{sweep\_selection.py} with the released simulator on the
same 15+32 training set; the released \texttt{sweep\_results.csv} records every
configuration, and the selected values match the frozen configuration replayed
by \texttt{run\_main\_study.sh}
(\texttt{--ncand 5 --kfrac 0 --z 1.28 --edelay 10}).

\subsection{Controller Under Test}

The controller (Fig.~1) keeps a fast path that serves every reference with the
active policy and updates two counters per page (last access and an aged
count); a sampled subset of references, selected by a 64-bit hash of the page
id with probability $p$, feeds reuse-time and window buffers. All other work
runs once per epoch of $W$ references. The candidate set is
\{LRU, LFU-A, 2Q, MRU, S3-FIFO\}, the five candidates selected on training traces
(Section~V-B); S3-FIFO is therefore inside DA's candidate set, but DA is not
significantly better than S3-FIFO in the reported comparisons. The initial
policy is LRU, and after the first baseline period one shadow evaluation
selects a start policy.

\begin{figure}[!t]
\centering
\includegraphics[width=0.85\columnwidth]{fig1_arch.png}
\caption{Architecture of the evaluated controller: per-reference fast path
(blue), per-epoch control path (orange), statistical decision points, and
policy-switching feedback.}
\label{fig:arch}
\end{figure}

\emph{Stage 1.} The epoch miss ratio $m_t$ is standardized as
$z_t = (m_t - \mu_0)/\sigma_0$, with $\mu_0$ and $\sigma_0$ estimated over a
baseline of $B = 8$ epochs ($\sigma_0$ floored at 0.02). A two-sided CUSUM
accumulates
\begin{align}
S_t^{+} &= \max(0,\ S_{t-1}^{+} + z_t - \delta) \label{eq:cusumplus}\\
S_t^{-} &= \max(0,\ S_{t-1}^{-} - z_t - \delta) \label{eq:cusumminus}
\end{align}
and raises an alarm when either statistic exceeds $h$.
\emph{Stage 2.} On an alarm, the sampled reuse times of the last three epochs
(measured in sampled references, with first-time references given a large
sentinel value) are compared with the baseline reuse times by a two-sample KS
test; drift is confirmed if $p < \alpha$ with at least 100 samples per side,
otherwise the alarm is rejected. Because reuse time depends on the workload and
not on the active policy, this stage is intended to separate workload change
from policy-induced miss-ratio change.

\emph{Shadow evaluation.} Ten epochs after confirmation (to accumulate
post-change data), each candidate is simulated at capacity
$C_s = \mathrm{round}(pC)$ on the sampled window starting at the confirmation
epoch. The first 40\% of the window warms the shadow caches, and the remainder
is split into four blocks. For candidate $\pi$, $\Delta_\pi$ is the per-block
miss reduction \emph{per sampled reference} relative to the active policy's
shadow, $B_{\mathrm{raw}}$ the raw references per block, $p$ the sampling rate,
$H$ a fixed horizon and $K = 0.1C$ a transition-cost margin in miss-equivalents.
Dividing by $p$ rescales sampled misses to raw references, so the controller
switches to the candidate with the largest margin if
\begin{equation}
\bigl(\mathrm{LCB}(\Delta_\pi) / (p\,B_{\mathrm{raw}})\bigr) H - K > 0
\label{eq:switch}
\end{equation}
where LCB is a one-sided lower confidence bound ($z = 1.28$) across the four
blocks. The frozen configuration sets $K = 0$, so the margin term did not
constrain the reported switches. A reverse switch to the previous policy within 30 epochs must
additionally exceed $K$ (hysteresis), and a dwell time of 10 epochs applies
after any switch. Switching rebuilds the target policy's metadata from the
resident pages ordered by last access, after discarding pages not referenced
in the last 20,000 references (A1); LFU-A and 2Q use the aged counts to place
pages. After a confirmed change the controller re-baselines for $4 + 10$ epochs
plus $B$ epochs. Table~V lists all parameters. The persistence safeguard of the
design (several consecutive confirmed epochs) was not implemented; the
monitored features were the miss ratio (Stage 1) and sampled reuse times
(Stage 2) only.

\begin{table}[!t]
\caption{Controller Parameters Fixed After Training}
\label{tab:params}
\footnotesize
\begin{tabular}{@{}p{0.42\columnwidth}p{0.52\columnwidth}@{}}
\toprule
Parameter & Value \\
\midrule
Capacity $C$; epoch $W$ & 1,000 frames; 2,000 references \\
Sampling rate $p$ & 0.2 (hash-based, spatial) \\
CUSUM $\delta$, $h$; baseline $B$ & 0.5, 5; 8 epochs ($\sigma_0 \ge 0.02$) \\
KS level $\alpha$; window & $10^{-3}$; last 3 epochs, $\ge$ 100 samples \\
Evaluation delay; cooldown & 10 epochs; 4 epochs \\
Shadow capacity; warm-up; blocks & $\mathrm{round}(pC) = 200$; 40\%; 4 \\
LCB $z$; horizon $H$; margin $K$ & 1.28; 50,000 refs; $K = 0$ in frozen run (design default $0.1C = 100$) \\
Dwell $D$; hysteresis window & 10 epochs; 30 epochs \\
Migration purge age & 20,000 references \\
Count aging & halved every 50 epochs \\
\bottomrule
\end{tabular}
\end{table}

\subsection{Metrics and Statistical Analysis}

The primary metric is the miss ratio after warm-up, reported with the
percentage reduction relative to LRU. Detection is evaluated on all 80
non-stable traces across four scenarios (abrupt, gradual, recurring and
composite; 380 true boundaries in total, which is what the 278/380 figure
denominates): a boundary is detected if a confirmed alarm occurs within
30,000 references after it, and the delay is the distance to that alarm. A
confirmed alarm outside every such window is counted as spurious. On stable
traces confirmations are counted both over the whole trace and after reference
100,000, the latter excluding an initial transient. A reversal is a switch
A$\to$B$\to$A within 60,000 references. Cost is the simulated wall-clock time per reference of the single
released simulator on this machine; it is indicative, not a kernel measurement.
Recovery time, memory overhead and full ranking agreement were not measured in
this reproducible study.

The unit of analysis is the trace (80 non-stable, 30 stable). For every
comparison in this paper the primary uncertainty statement is the percentile
interval over traces (10,000 resamples for means, 5,000 for paired
differences); paired differences are also tested with the Wilcoxon signed-rank
test [33] and Holm correction [35] across the ten baseline comparisons, and
effect sizes are mean and median paired differences with win/loss counts.
Traces in one scenario share a generator, so they are not independent draws
from a population of real workloads; only five access primitives and five
kernel models underlie all traces, and the effective number of independent
units is smaller than 80. The trace-level intervals are therefore conditional
on this design. \emph{Sensitivity only.} As a coarse check we also resample the
four scenario clusters; with four clusters this yields only $\binom{7}{4}=35$
distinct resamples. It gives DA--S3-FIFO $-0.0154$ [$-0.0294$, $-0.0048$],
excluding zero where the primary trace-level paired interval does not. Because
35 resamples cannot support a reliable interval, this result is reported only
as a sensitivity note and is never used as the headline uncertainty statement.



\subsection{Replacement Quality}

All numbers in this section come from the frozen configuration and the fresh
test traces; the raw per-trace CSV output is in the reproduction package.
Table~VI and Figs.~3--5 give the primary result. On the 80 non-stable traces
the controller's mean miss ratio was 0.453 (95\% CI 0.423--0.484), 20.5\%
below LRU (0.570). The offline optimum reached 0.302. The controller was
better than S3-FIFO (0.468) on the mean, but the paired difference was not
significant (Holm-adjusted $p = 0.09$; Table~VII). It was significantly
\emph{worse} than continuous shadow evaluation (ContShadow, 0.423) by 0.030
miss ratio (adjusted $p < 0.001$). On the 30 stable traces the controller
scored 0.307, close to ContShadow (0.307) and better than S3-FIFO (0.330) and
LRU (0.519), largely because of start-up policy selection.

\emph{Stable-trace gain is not drift-handling gain.} A stable trace contains no
phase change, so a benefit there cannot be attributed to drift detection. On
the 30 stable traces DA scored 0.3071 with 0.80 switches and 1.73 shadow
evaluations per trace, against LRU 0.5186, S3-FIFO 0.3301 and ContShadow
0.3065. That advantage over LRU comes from the initial shadow evaluation that
chooses a start policy and from continued re-evaluation, not from responding
to drift. The drift-handling comparison is the non-stable result above, where
DA is 0.030 worse than ContShadow despite saving 76\% of shadow work. We do
not add the stable-trace gain to the drift-handling claim.

By scenario the controller was close to ContShadow on abrupt (0.370 versus
0.365) and recurring (0.398 versus 0.384) traces, but clearly behind it on
gradual drift (0.439 versus 0.401) and composite (0.604 versus 0.541). The
gap on gradual drift matters because the paper's premise is detection of slow
change; the CUSUM trigger usually reacts only after the mixture has moved.

\begin{table}[!t]
\caption{Mean Miss Ratio by Scenario and Overall (Fresh Test Traces)}
\label{tab:quality}
\footnotesize
\setlength{\tabcolsep}{2pt}
\begin{tabular}{@{}lccccc@{}}
\toprule
Method & Abrupt & Gradual & Recurring & Composite & Non-stable \\
\midrule
LRU & 0.500 & 0.572 & 0.527 & 0.682 & 0.570 [0.540, 0.598] \\
FIFO & 0.536 & 0.609 & 0.560 & 0.696 & 0.600 [0.573, 0.627] \\
CLOCK & 0.487 & 0.558 & 0.515 & 0.677 & 0.559 [0.530, 0.589] \\
LFU-A & 0.474 & 0.535 & 0.505 & 0.671 & 0.546 [0.515, 0.576] \\
MRU & 0.943 & 0.928 & 0.902 & 0.968 & 0.935 [0.918, 0.949] \\
2Q & 0.457 & 0.504 & 0.479 & 0.648 & 0.522 [0.493, 0.551] \\
ARC & 0.465 & 0.520 & 0.496 & 0.634 & 0.529 [0.499, 0.557] \\
S3-FIFO & 0.406 & 0.441 & 0.412 & 0.615 & 0.468 [0.435, 0.501] \\
LeCaR & 0.476 & 0.545 & 0.503 & 0.680 & 0.551 [0.520, 0.582] \\
LPR-style & 0.477 & 0.535 & 0.474 & 0.646 & 0.533 [0.509, 0.556] \\
ContShadow & 0.365 & 0.401 & 0.384 & 0.541 & \textbf{0.423} [0.400, 0.446] \\
DA & 0.370 & 0.439 & 0.398 & 0.604 & \textbf{0.453} [0.423, 0.484] \\
OPT & 0.244 & 0.269 & 0.242 & 0.452 & 0.302 [0.281, 0.324] \\
\bottomrule
\end{tabular}
\end{table}

\begin{figure}[!t]
\centering
\includegraphics[width=0.85\columnwidth]{fig3_overall.png}
\caption{Mean miss ratio on the 80 non-stable traces with 95\% bootstrap
intervals (reproducible main study, $C = 1{,}000$).}
\label{fig:overall}
\end{figure}

\begin{figure}[!t]
\centering
\includegraphics[width=0.85\columnwidth]{fig4_scenario.png}
\caption{Mean miss ratio by scenario on the fresh test traces.}
\label{fig:scenario}
\end{figure}

\begin{table}[!t]
\caption{Paired Differences, DA Minus Baseline (Non-Stable, $N = 80$)}
\label{tab:paired}
\footnotesize
\setlength{\tabcolsep}{4.5pt}
\begin{tabular}{@{}lcccc@{}}
\toprule
Baseline & Mean diff. [95\% CI] & W/L/T & Holm $p$ \\
\midrule
LRU & $-0.1171$ [$-0.140$, $-0.094$] & 75/5/0 & $<0.001$ \\
FIFO & $-0.1472$ [$-0.171$, $-0.124$] & 75/5/0 & $<0.001$ \\
CLOCK & $-0.1065$ [$-0.130$, $-0.084$] & 75/5/0 & $<0.001$ \\
LFU-A & $-0.0933$ [$-0.116$, $-0.071$] & 74/6/0 & $<0.001$ \\
2Q & $-0.0691$ [$-0.090$, $-0.048$] & 71/9/0 & $<0.001$ \\
ARC & $-0.0758$ [$-0.099$, $-0.053$] & 71/9/0 & $<0.001$ \\
S3-FIFO & $-0.0154$ [$-0.035$, $+0.004$] & 42/38/0 & 0.090 \\
LeCaR & $-0.0980$ [$-0.121$, $-0.076$] & 72/8/0 & $<0.001$ \\
LPR-style & $-0.0800$ [$-0.098$, $-0.061$] & 70/10/0 & $<0.001$ \\
ContShadow & $+0.0301$ [$+0.016$, $+0.045$] & 15/54/11 & $<0.001$ \\
\bottomrule
\end{tabular}
\end{table}

Table~VII reports paired comparisons; negative differences favour the
controller. It was significantly better than eight of the ten baselines after
Holm correction (all adjusted $p < 0.001$). The exception was S3-FIFO, where
the mean favoured DA by only 0.015 and the adjusted $p$ was 0.09: the evidence
for an improvement over the strongest fixed policy is therefore weak. The
other exception is unfavourable: ContShadow had a 0.030 lower mean miss ratio
(adjusted $p < 0.001$), so continuous shadow evaluation was more accurate than
drift-gated evaluation in this study. Fig.~5 shows the paired differences.
Bootstrap intervals describe uncertainty in the mean; the Wilcoxon test uses
ranks, so the two need not agree. Intervals are not corrected for multiplicity
or for clustering of traces by scenario.

\emph{Post hoc design correction (synthetic traces).} The capacity-sampling
repair described in Section~\ref{sec:public} was re-selected on the training
traces only using the same 48-configuration grid and selection rule with
\texttt{--cs-cap-f 1.0}. It selected the same configuration
($n_{\mathrm{cand}}{=}5$, $K{=}0$, delay 10). Evaluated once on the 80
non-stable and 30 stable held-out test traces, the repaired controller
produced miss ratios, switch counts, and shadow-evaluation counts identical to
the original controller. This post hoc correction therefore did
not change the synthetic-trace outcome.

\begin{figure}[!t]
\centering
\includegraphics[width=0.85\columnwidth]{fig5_paired.png}
\caption{Paired differences in miss ratio between DA and each baseline on the
80 non-stable traces (negative favours DA).}
\label{fig:paired}
\end{figure}

\subsection{Ablation}

Table~\ref{tab:ablation} removes one component at a time. The dominant result
is that shadow evaluation is essential: replacing it by a random choice raised
the non-stable miss ratio from 0.453 to 0.619, worse than LRU. The remaining
variants scored 0.453--0.456, a small range that must be interpreted with the
activation counts below. Additive diagnostic logging of the frozen main-study
run (byte-identical outputs to the released simulator) records how often each
component actually acted across the 110 traces. The recency filter acted at
47 of the 213 switches, dropping 39,405 stale pages from the migrated resident
set. The KS stage was active on all 643 alarms and rejected 173 of them, after
which 570 shadow evaluations were performed; its small ablation effect is not
because it never fired. The cost margin never acted ($K = 0$, so no cost-margin
decision was made) and the dwell rule never blocked a switch (zero blocks):
under the frozen configuration these two components were \emph{inactive}, and
the ablation cannot estimate their benefit from these runs. Counting an
inactive component as ``showing no benefit'' would conflate absence of
activation with absence of effect. Table~\ref{tab:ablation} also includes
S3-FIFO as a static reference: at 0.468 miss ratio and 28.5 ns/reference it is
cheaper than DA but less accurate, so neither policy Pareto-dominates the
other on these generated traces. ContShadow reached 0.423 with 27.39 shadow
evaluations per trace against 6.47 for DA, a 76.4\% reduction, but it was
more accurate: the saving in shadow work was \emph{not} at equal miss ratio on
the generated traces. On the real-program traces below the two methods do tie
on miss ratio.

\begin{table}[!t]
\caption{Ablation on the Fresh Test Traces (Non-Stable Mean)}
\label{tab:ablation}
\footnotesize
\setlength{\tabcolsep}{5pt}
\begin{tabular}{@{}lcccc@{}}
\toprule
Variant & Miss ratio & Switches & Shadow evals & ns/ref \\
\midrule
DA & 0.453 & 2.36 & 6.47 & 67.4 \\
DA-noPurge & 0.453 & 2.36 & 6.47 & 68.5 \\
DA-noS2 & 0.456 & 2.27 & 6.76 & 67.9 \\
DA-noShadow & 0.619 & 5.05 & 0.00 & 77.3 \\
DA-noCost & 0.453 & 2.36 & 6.47 & 66.9 \\
DA-noDwell & 0.453 & 2.36 & 6.47 & 67.9 \\
S3-FIFO (static reference) & 0.468 & 0.00 & 0.00 & 28.5 \\
ContShadow & 0.423 & 3.15 & 27.39 & 133.4 \\
\bottomrule
\end{tabular}
\end{table}

\begin{figure}[!t]
\centering
\includegraphics[width=0.85\columnwidth]{fig6_ablation.png}
\caption{Ablation: non-stable miss ratio and shadow evaluations per trace.}
\label{fig:ablation}
\end{figure}

\subsection{Detection, Switching and Cost}

The CUSUM trigger detected 278 of 380 phase boundaries (73.2\%) over the four
non-stable scenarios, with a median confirmation delay of 2,000 references:
97/100 abrupt, 24/100 gradual, 96/100 recurring and 61/80 composite. Detection
is therefore strong for abrupt and recurring changes and weak for gradual
ones. This value is conditional on the boundary alignment described below.
Boundaries in the generator coincide with epoch boundaries, which favours
detection, so neither figure should be read as a real-system detection rate.
Of the 169 spurious confirmations outside the 30,000-reference boundary
windows, 16 occurred on abrupt, 109 on gradual, 14 on recurring and 30 on
composite traces: the gradual scenario, already the hardest to detect,
produced about two-thirds of all false alarms. On the 30 stable traces there
were 23 spurious confirmations, but only one after reference 100,000. Of the
23, 17 occurred before reference 50,000, five between 50,000 and 100,000, and
the remaining one at reference 160,000; the 100,000-reference exclusion thus
removes 22 of the 23, and the residual late confirmation is a genuine
long-horizon false alarm rather than an initial transient. The controller
switched 2.36 times per non-stable trace against 3.15 for ContShadow.

\emph{Detection with shifted phase boundaries.} The generator's change points
fall on controller-epoch boundaries, which favours detection. The released
simulator's \texttt{--phase-offset} option prepends a given number of
references from the first phase before the first boundary and shifts all
recorded boundaries by that amount; with offset zero the main detection output
is reproduced exactly. Table~S.VII (supplementary material) reports seven offsets generated
directly from raw detection CSVs by \texttt{generate\_offset\_table.py},
including 333 and 1501 references, which are not multiples of the controller
epoch size ($W{=}2000$). Detection is counted within a 30,000-reference
absolute window after each recorded boundary. Across offsets the trigger
detected between 274 and 281 of 380 boundaries (72.1\%--73.9\%) and produced
between 169 and 173 spurious confirmations; offset 0 was neither best nor
worst. Stable traces produced 23 spurious confirmations at every offset.

Gradual drift remains weakly detected: 24 of 100 gradual boundaries were
detected at offset 0 (17--26 across offsets), while gradual traces produced 109
spurious confirmations at offset 0. This means the paper's drift-gating premise
is supported only for sharper regime changes in this study; slowly moving
workloads can cross the miss-ratio threshold without a reliable CUSUM
confirmation.



The raw CSV rows, the exact simulator commands and the diagnostic event-log
schema are collected in the supplementary material accompanying this paper.

Simulated time per reference was 25.1 ns for LRU, 8.9 for CLOCK, 28.5 for
S3-FIFO, 36.8 for 2Q, 50.2 for LPR-style, 67.4 for DA and 133.4 for
ContShadow: DA cost 2.7 times LRU and 2.4 times S3-FIFO. These are relative
figures from an unoptimized single-threaded simulator, not kernel latency. DA
is slower and, in this study, less accurate than ContShadow, but cheaper;
against S3-FIFO it is both slower and (non-significantly) more accurate. No
configuration here beat the best static policy on both cost and quality.

\begin{figure}[!t]
\centering
\includegraphics[width=0.85\columnwidth]{fig7_cost.png}
\caption{Miss ratio against simulated time per reference (indicative,
single-threaded simulator); S3-FIFO is included as a static comparator.}
\label{fig:cost}
\end{figure}

\subsection{Post Hoc Measurement Experiments}
\label{sec:posthoc}
The measurements below were run after the frozen main-study results were
obtained. They are reported as post hoc measurements, not as primary
claims or as selected configurations. Raw outputs are in
\texttt{reproduction/trigger\_real/}, \texttt{reproduction/param\_sweep/},
and \texttt{reproduction/rankability\_real/}; commands, seeds, and defaults
are recorded in \texttt{NOT\_DONE.md}.

\emph{Trigger variants (synthetic, post hoc).} The released simulator was
modified to add \texttt{--trigger cusum|periodic|random|oracle},
\texttt{--trigger-period}, \texttt{--trigger-rand}, \texttt{--trigger-seed},
and \texttt{--delta}; default CUSUM DA output was verified byte-identical to
the released main-study DA rows (110 rows, 0 mismatches, detection
identical). On the same 80 non-stable seed-5001 test traces, DA-oracle
(mean miss 0.4335, 2.19 switches, 5.74 shadow evals) was better than DA
(0.4529, 2.36, 6.47), while DA-periodic-p40 (0.4687, 2.52, 7.22) and
DA-random-r04 (0.4906, 2.60, 6.50) were worse. Paired DA-minus-other
differences (bootstrap 10,000, seed 0; Wilcoxon/Holm) were: DA-oracle
$+0.0194$ [$+0.0061,+0.0347$], W/L/T 8/22/50, Holm $0.0177$; DA-periodic-p40
$-0.0158$ [$-0.0297,-0.0008$], W/L/T 36/26/18, Holm $0.0177$; DA-random-r04
$-0.0377$ [$-0.0585,-0.0163$], W/L/T 47/19/14, Holm $0.0004$; DA-ContShadow
$+0.0301$ [$+0.0167,+0.0454$], W/L/T 15/54/11, Holm $2.51{\times}10^{-5}$.
The oracle variant uses synthetic ground-truth boundaries and is not a
deployable trigger.

\emph{FIU trigger check (post hoc).} Using the same six FIU test traces and
base args as the FIU rankability runs (default seed 5001, default $C=1000$),
DA-oracle had mean miss 0.712698, 6 switches, 6 shadow evals, and 0 confirms;
DA-periodic-p40 had 0.713068, 11, 40, 35 confirms; DA-random-r04 had
0.709742, 10, 36, 35 confirms. The corresponding baseline at $C=1000, p=0.2$
was 0.706285 mean miss, 10 switches, 46 shadow evals
(\texttt{rankability\_real/fiu\_C1000\_p0.2.csv}). The oracle produced no
confirmations because these FIU trace files expose no ground-truth bounds in
the trace format; this result is not evidence of oracle benefit on FIU.

\emph{Parameter sensitivity (synthetic, post hoc).} One-at-a-time sweeps from
the default DA configuration (seed 5001, 80 non-stable + 30 stable traces)
showed: $W\in\{1000,2000,4000\}$ gave miss 0.461442, 0.452888, 0.500050;
$\delta\in\{0.1,0.3,0.5,0.7,1.0\}$ gave 0.457896, 0.455087, 0.452888,
0.454430, 0.453678; $h\in\{2,3,5,7,10\}$ gave 0.456338, 0.453806, 0.452888,
0.454700, 0.456900; $p\in\{0.05,0.1,0.2,0.4\}$ gave 0.480957, 0.469601,
0.452888, 0.447180. These measurements describe local sensitivity around the
frozen configuration and do not replace the train-only selection protocol.

\emph{Capacity/sampling rankability (post hoc).} The synthetic sweep
($C\in\{500,1000,2000,4000\}$, $p\in\{0.05,0.1,0.2,0.4\}$) showed mean miss
decreasing with larger $C$ (e.g., $C=500,p=0.2$ 0.580196; $C=4000,p=0.2$
0.180884) with nonzero switches/evals in all tested cells
(\texttt{rankability\_real/synth\_sweep.csv}). The FIU sweep
($C\in\{500,1000,2000,5000,10000,20000\}$, same $p$ grid) showed
$C=20000$ gives mean miss 0.668337 with 0 switches for all $p$, while smaller
$C$ values show switches and miss variation
(\texttt{rankability\_real/fiu\_sweep.csv}). These real measurements
characterize rankability/capacity behavior but do not provide a new external
validation win.

\subsection{Public Block-Trace Check}
\label{sec:public}
The generated traces have controlled change points, which favours the
controller. As an external check the released simulator was run on six FIU
web/block traces from the CMU PDL cacheDatasets collection [39], obtained from
the repository's \texttt{fiu/} directory. The released converter
\texttt{convert\_public\_traces.py} reads the packed 24-byte request records
(timestamp, object id, object size, next-access time) and maps object ids to
dense identifiers in first-touch order. The first 1,000,000 requests of each
trace were used, except \texttt{madmax}, which contains 918,010 requests; the
single-thread page-id array requires \texttt{--np 1048576}. The replay
capacity $C = 20{,}000$ was selected on a \emph{separate} PDL trace,
\texttt{fiu\_webusers-030409-033109.oracleGeneral.zst}, which is not one of the
six evaluated traces. The recorded post hoc grid was $C \in \{1000, 5000, 10000,
15000, 20000, 25000, 50000\}$ with seed 0, methods LRU and S3-FIFO (DA was also
run for the record). The rule selected the smallest $C$ at which both miss
ratios were at least 0.01 and LRU exceeded S3-FIFO by at least 0.10: at
$C = 10{,}000$ the miss ratios were 0.74182/0.73317 (gap 0.00865), and at
$C = 20{,}000$ they were 0.73590/0.52726 (gap 0.20864), so $C = 20{,}000$ was
selected. The raw pilot CSVs, manifest and exact commands are in
\texttt{phase10\_capacity\_pilot/}. The earlier \texttt{webresearch}-only pilot
(\texttt{public\_pilot\_C*.csv}) is retained for provenance but is not the
selection source. The frozen main-study controller configuration and the
default seed 5001 apply to the six evaluated traces. The exact replay commands
are given in the supplementary material.

\begin{table}[!t]
\caption{Public FIU Block Traces: Means Over Six Traces (Raw Output
\texttt{public\_trace\_results.csv})}
\label{tab:public}
\footnotesize
\setlength{\tabcolsep}{4pt}
\begin{tabular}{@{}lccc@{}}
\toprule
Method & Mean miss ratio & Mean ns/ref & Mean shadow evals \\
\midrule
OPT & 0.4595 & 0.0 & 0.00 \\
S3-FIFO & 0.5772 & 31.7 & 0.00 \\
2Q & 0.6209 & 38.3 & 0.00 \\
ARC & 0.6580 & 37.0 & 0.00 \\
LRU & 0.6753 & 25.5 & 0.00 \\
DA & 0.6753 & 82.5 & 13.50 \\
ContShadow & 0.6753 & 180.1 & 48.33 \\
FIFO & 0.6808 & 17.6 & 0.00 \\
\bottomrule
\end{tabular}
\end{table}

Table~\ref{tab:public} reports the means. DA, ContShadow and LRU have the same
miss ratio on every one of the six traces (mean 0.6753), and DA made zero
switches on every trace, with 11--15 KS-confirmed alarms per trace and no true
boundaries. S3-FIFO had both a lower mean miss ratio (0.5772) and a lower mean
simulated time (31.7 versus 82.5 ns/reference), so DA is \emph{Pareto-dominated}
by S3-FIFO on this trace set; DA is also weakly dominated by LRU, which matches
its miss ratio at 25.5 ns/reference. Gating saved 72\% of the shadow
evaluations of continuous evaluation here but yielded no accuracy gain.

\emph{Per-trace diagnostics and a capacity--sampling fault.} To check for a
sampling or implementation fault, each public trace was re-run with an
additive diagnostic build of the released simulator that logs the CUSUM
statistic and per-candidate shadow margins without changing control flow
(identical miss ratios, switch counts and shadow-evaluation counts; byte-identical
detection CSV). The CUSUM alarm threshold $h = 5$ was crossed on every trace:
maximum $S^+$ ranged from 7.95 to 20.30 and maximum $S^-$ from 6.08 to 33.30,
so the trigger fired normally. Yet every shadow evaluation returned identical
miss counts for all five candidates --- for example, 1,120 misses for LRU,
LFU-A, 2Q, MRU and S3-FIFO on \texttt{webresearch} at reference 16,000 --- giving
zero per-candidate margins and a tie-break to candidate 0 (LRU) in every
initial selection and every later decision. S3-FIFO is a candidate in DA
(candidate id 4), so its 0.577 stand-alone miss ratio was simply never visible
to the controller. The cause is an interaction between capacity and sampling:
the shadow cache holds $C_s = p\,C = 0.2 \times 20{,}000 = 4{,}000$ pages, more
than the roughly 1,100--1,800 distinct pages in each sampled evaluation window,
so all candidates incur only compulsory misses and cannot be ranked. At the
synthetic-study capacity $C = 1{,}000$ ($C_s = 200$) the same trace gives
differentiated counts (1,141 / 1,137 / 1,139 / 1,885 / 1,135 for candidates
0--4 at reference 16,000), S3-FIFO wins the initial selection and DA makes one
switch. The no-switch result on the public traces is therefore a
capacity--sampling fault in this configuration, not evidence that candidate
policies are equivalent on real workloads. Per-trace maxima and margins are in
the supplementary material.

\emph{Parameter re-selection on training prefixes.} Each public trace was split
into a 500,000-reference training prefix and a held-out suffix (418,010
references for \texttt{madmax}; 500,000 for the others), and the 48-configuration
selection grid was re-run on the six prefixes. All 48 configurations produced
the same training mean miss ratio (0.68027) and zero switches, because the
capacity--sampling fault makes the controller parameters inert at $C = 20{,}000$;
the tie-break selected $n_{\mathrm{cand}} = 4$, $K = 0$, evaluation delay 2. On
the held-out suffixes the re-selected configuration and the frozen one
($n_{\mathrm{cand}} = 5$, $K = 0$, delay 10) both achieved mean miss 0.66834 with
zero switches; the re-selected run was cheaper per reference (67.6 versus 85.0
ns) only because it evaluated four candidates instead of five. These runs are
reported as a failed external validation, not as a tuned result.

\emph{Post hoc design correction.} A post hoc repair caps shadow capacity at
$C_s=\min(\mathrm{round}(pC),\mathrm{round}(f D_w))$, where $D_w$ is the
number of distinct sampled page IDs in the evaluation window and $f=1.0$ is
fixed a priori. With $f=1.0$ this is an identity cap: it changes $C_s$ only
when $D_w < pC$, but when all candidates fit in the sampled distinct-page
set they still tie on compulsory misses. It therefore cannot create
rankability where $pC$ already exceeds $D_w$; a real fix must reduce $C_s$
below $D_w$ or increase $D_w$ relative to $C_s$. The repair is implemented
behind \texttt{--cs-cap-f} (default 0, disabled) and leaves the original
behaviour byte-identical. Re-running the 48-configuration selection on the
same training prefixes with \texttt{--cs-cap-f 1.0} produced the same training
mean miss ratio (0.68027) and zero switches for every configuration; the
tie-break selected the same configuration. On the held-out suffixes, the
repaired controller achieved the same mean miss ratio as the original
(0.66834) with zero switches. The post hoc repair therefore did not change
the public-trace outcome.

These are storage block/cache traces, not virtual-memory page traces, and only
a prefix of each was replayed; this is an external check and not the primary
evidence for paging.

\subsection{Post Hoc Follow-up Measurements}
\label{sec:posthoc_followups}

The following measurements are recorded in \texttt{PROTOCOL.md}
(first committed to git on 2026-11-24; no independent pre-run timestamp
exists) and are reported as post hoc. They test repairs and
alternatives suggested by the negative main result; null and negative
outcomes are reported.

\begin{table}[!t]
\caption{Post hoc follow-ups: outcome and representative numbers. Details
and raw-table pointers are in Section S4 of the supplementary material.}
\label{tab:prespec}
\footnotesize
\setlength{\tabcolsep}{3pt}
\begin{tabular}{@{}p{0.18\columnwidth}p{0.74\columnwidth}@{}}
\toprule
Follow-up & Outcome (key numbers) \\
\midrule
Phase 1 rankability & $C_s/D_w\ge1$ gives 0/1514 rankable evals; no FIU repair gain (C=20000 test 0.6621). \\
Phase 2 budget frontier & CUSUM better than periodic/random near 6.2--7.0 evals; Oracle 0.4335, ContShadow 0.4228 accuracy ceiling. \\
Phase 3 detectors & PH 0.4552, ADWIN 0.4991 vs CUSUM 0.4529; neither closes oracle gap. \\
Phase 4 baselines/traces & SIEVE 0.5557, W-TinyLFU 0.4913, LHD 0.6736 synth; real-trace DA no advantage; PDL FIU originals: S3-FIFO 0.6188 best, DA 0.6280. \\
Phase 5 robustness & 2.2\,$\mu$s DA/ContShadow crossover is calibration only; t3 LCB DA 0.4586 vs 0.4529 (negative); DA never on total-time envelope. \\
\bottomrule
\end{tabular}
\end{table}

\subsection{Cost Break-Even and Total-Time Envelope}
\label{sec:breakeven}

The measured trade-off in Section~V-F is a 0.030 miss-ratio penalty for DA
against a 66.0 ns/reference time saving and 20.9 fewer shadow evaluations per
trace. Whether that is a good trade depends on what a miss costs. Let $c_m$ be
the cost of one miss (page fault plus service) and $c_e$ the cost of one shadow
evaluation, both in nanoseconds. With the same fast path for both methods, the
total time per trace is
\begin{equation}
T = N t_0 + N m c_m + E c_e,
\label{eq:cost}
\end{equation}
where $N$ is the trace length, $m$ the miss ratio, $E$ the number of shadow
evaluations, and $t_0$ the per-reference cost excluding shadows. DA beats
ContShadow in total time iff
\begin{equation}
(m_{\mathrm{DA}} - m_{\mathrm{CS}})\,c_m
+ (E_{\mathrm{DA}} - E_{\mathrm{CS}})\,c_e/N < 0,
\label{eq:breakeven}
\end{equation}
i.e.\ $c_e > R c_m$ with
$R = N(m_{\mathrm{DA}} - m_{\mathrm{CS}})/(E_{\mathrm{CS}} - E_{\mathrm{DA}})$.
Using the measured means ($m_{\mathrm{DA}}=0.453$, $m_{\mathrm{CS}}=0.423$,
$E_{\mathrm{DA}}=6.47$, $E_{\mathrm{CS}}=27.39$, and $N=545{,}000$ for the
average non-stable trace), \texttt{cost\_model.py} recomputes $R = 783.28$,
i.e.\ $R \approx 783$: one shadow evaluation must cost more than about 783
misses for DA to beat ContShadow. A 10,000-replicate trace-level bootstrap
(\texttt{bootstrap\_cost\_model.py}, seed 20261123) gives a 95\% percentile CI
of $[426, 1187]$ for $R$. Equivalently, the direct measured-time comparison
$T \approx t + m c_m$ gives a DA--ContShadow crossover of
$(133.4-67.4)/(0.453-0.423) = 2196\,\mathrm{ns} \approx 2.2\,\mu$s per miss
(\texttt{cost\_model.csv}), with a 95\% bootstrap CI of
$[1.44, 4.06]\,\mu$s (\texttt{cost\_model\_bootstrap.csv}). Both figures depend
on the unoptimized simulator timings; the ratio is more robust than the
absolute scale, and the intervals are wide because the 80-trace sample has a
heavy tail. These are DA-versus-ContShadow comparisons; Table~\ref{tab:breakeven}
gives the full total-time comparison over four policies and
Fig.~\ref{fig:breakeven} plots it. S3-FIFO is already a candidate in DA's shadow
set (candidate id 4 in \texttt{mkCand}); DA pays the gating and evaluation
overhead even when it selects S3-FIFO.

\textbf{Total-time envelope.}
Using the measured means (CLOCK 0.559 miss ratio at 8.90\,ns/reference,
S3-FIFO 0.468 at 28.54\,ns, ContShadow 0.423 at 133.36\,ns, DA 0.453 at
67.36\,ns), the lower envelope is CLOCK for $c_m<0.216\,\mu$s, S3-FIFO for
$0.216<c_m<2.305\,\mu$s and ContShadow for $c_m>2.305\,\mu$s. DA is never on
the envelope: its smallest excess over the best policy is
3.27\,ns/reference at $c_m\approx2.305\,\mu$s. At the measured minor-fault
cost (1.98--2.06\,$\mu$s) S3-FIFO is cheapest and DA is about 5\,ns/reference
above it. DA beats ContShadow only below 2.196\,$\mu$s, but S3-FIFO stays
cheaper than DA until 2.517\,$\mu$s. A 10,000-replicate trace-level bootstrap
(seed 20261124, \texttt{phase5/bootstrap\_total\_time.py}) gives 95\% intervals
of $[0.171,0.274]\,\mu$s (CLOCK/S3-FIFO), $[1.660,3.678]\,\mu$s
(S3-FIFO/ContShadow), $[1.449,4.105]\,\mu$s (DA/ContShadow) and
$[1.086,22.05]\,\mu$s (DA/S3-FIFO; unstable because the miss-ratio difference
is small); the minimum-gap interval is $[-28.6,40.2]$\,ns, so the ranking at
the crossing is not statistically resolved. The point estimates nevertheless
show no miss cost at which DA is cheapest.

To check the implied evaluation cost directly, a timing-only build of the
released simulator (no diagnostic logging) measured 570 shadow evaluations on
the same traces at a mean of $2.46$\,ms and a median of $2.33$\,ms each. The
cost model's time decomposition implies $1.72$\,ms per evaluation, so the
direct measurement is about $1.4\times$ higher; both are simulator-side
numbers from an unoptimized single-threaded build and should be read as an
order-of-magnitude check, not a kernel latency. Because the directly measured
evaluation cost exceeds the implied one, the DA--ContShadow crossover would
occur at a somewhat lower miss cost than the $2.2\,\mu$s point estimate, which
further weakens the case for gating under minor-fault-dominated costs.

\input{total_time_table.tex}

The 2.2\,$\mu$s point is a user-space minor-fault calibration
(\texttt{prototype\_fault\_latency.cpp}), not a kernel page-fault latency.
On the evaluation machine, a first-touch minor fault on a 4\,KiB anonymous page
costs 1.98\,$\mu$s and a refault after \texttt{MADV\_DONTNEED} costs
2.06\,$\mu$s (medians over seven 128\,MiB runs; a resident access costs
16\,ns). The other $c_m$ rows are illustrative parameter ranges, not measured
on this system; they show how the DA--ContShadow crossover moves with miss
penalty. At this calibration S3-FIFO is cheapest and DA is close to
ContShadow; for larger miss penalties ContShadow's lower miss ratio dominates.
Across all miss costs DA is above the lower envelope of CLOCK, S3-FIFO and
ContShadow (Table~\ref{tab:breakeven}). An analytical async-shadow variant
(\texttt{phase5/async\_model.csv}, assumption: evaluation fully off the
request path, same miss ratios, background CPU proportional to evaluations)
removes evaluation latency from the critical path but leaves the miss-ratio
difference unchanged; it therefore does not convert DA into an accuracy win
over ContShadow. Absolute times are from an unoptimized single-threaded
simulator; the DA--ContShadow ratio is the more robust quantity.

\begin{figure}[!t]
\centering
\includegraphics[width=0.85\columnwidth]{fig8_breakeven.png}
\caption{Total time per reference versus miss cost for CLOCK, S3-FIFO,
ContShadow and DA, with the lower envelope of the three comparators (black
dashed). CLOCK is cheapest below 0.216\,$\mu$s, S3-FIFO between 0.216 and
2.305\,$\mu$s, and ContShadow above 2.305\,$\mu$s; DA is never on the envelope,
with a minimum excess of 3.27\,ns/reference at 2.305\,$\mu$s. The orange band
is the measured minor-fault cost; the inset zooms the CLOCK/S3-FIFO crossing.}
\label{fig:breakeven}
\end{figure}

\subsection{Capacity Sensitivity}

Table~S.VIII (supplementary material) repeats the primary comparison at $C = 500$. The controller scored
0.580 against LRU 0.693 and S3-FIFO 0.590, but was again behind ContShadow
(0.555). The $C = 2{,}000$ and sampling-rate runs are present in the released
raw CSVs; full-shadow ranking agreement was not measured, so we make no claim
from it.



\subsection{Hypothesis Outcomes}

We frame three hypotheses. H1: the drift-aware controller attains a lower miss
ratio than the strongest fixed baseline (S3-FIFO). H2: a two-stage gate (CUSUM
trigger plus KS confirmation) improves miss ratio and detection over a single
CUSUM trigger with the same sampled shadow evaluation; the frozen run used the
design margin $K = 0$. H3: the controller's simulated per-reference cost remains
competitive with fixed and adaptive baselines. A previously listed hypothesis
about sampled shadow evaluation ranking candidates consistently with
full-capacity evaluation was dropped because full ranking agreement was not
measured; no claim is made from it.
Table~\ref{tab:hypotheses} reports the outcomes. No cost budget or agreement
threshold was specified before the study; post hoc criteria are not treated as
prospectively registered results.

\begin{table}[!t]
\caption{Hypothesis Outcomes (Reproducible Main Study)}
\label{tab:hypotheses}
\footnotesize
\begin{tabular}{@{}p{0.06\columnwidth}p{0.30\columnwidth}p{0.56\columnwidth}@{}}
\toprule
 & Outcome & Evidence \\
\midrule
H1 & Not supported & DA 0.453 vs S3-FIFO 0.468 ($-0.015$, adjusted $p = 0.09$); better than LRU and seven other baselines. \\
H2 & Not supported & ContShadow had lower miss ratio; the KS stage was active on all 643 alarms but had a small miss-ratio effect (0.456 without KS vs.\ 0.453 with KS); detection rate 73.2\%. \\
H3 & Not met & $2.7\times$ LRU per reference; DA is never on the lower total-time envelope (minimum excess 3.27\,ns/ref at $c_m\approx2.30\,\mu$s). \\
\bottomrule
\end{tabular}
\end{table}

The design lessons are narrow. First, shadow
evaluation is the component that matters: without it the controller is worse
than LRU. Second, the CUSUM trigger with sampled shadow evaluation reduced
shadow work by 76\% relative to continuous evaluation under the frozen cost-margin
rule ($K = 0$), but here it cost 0.030 in mean miss ratio; gating is a
cost-saving mechanism, not an accuracy advantage in the generated study. Third,
the KS stage had a small miss-ratio effect, while the cost margin and dwell
rule were inactive in the frozen run ($K = 0$, zero dwell blocks) and the
recency filter acted at 47 of 213 switches. KS was enabled in the frozen run,
while the margin was $K = 0$; the
effective mechanism was therefore a CUSUM trigger with sampled shadow
evaluation, not a validated two-stage gate. Fourth, on the fifteen-program
appendix the gated controller and continuous evaluation tie on miss ratio, but
DA costs 2.4 times S3-FIFO and does not beat it. Fifth, the cost model shows
that DA beats ContShadow only when a miss costs less than about 2.2 $\mu$s,
but the full total-time envelope puts CLOCK below 0.216\,$\mu$s, S3-FIFO
between 0.216 and 2.305\,$\mu$s and ContShadow above 2.305\,$\mu$s, with DA
never cheapest on the measured means. We therefore report where gating saves
shadow work, and do not claim that DA is best in these measurements or superior
to strong replacement policies.

\section{Threats to Validity}
\label{sec:threats}

\emph{Construct validity.} The traces are generated: five synthetic access
primitives and five kernel-level access models, not traces recorded from
running systems. Generated phases have clean boundaries aligned with the controller's epochs,
which favours detection delay, and workload types such as loops larger than
memory are overrepresented relative to their frequency in practice, which
favours policies such as MRU and regret-driven LRU/MRU switching. The reported
73.2\% detection rate (278/380) and roughly 2,000-reference
median delay are therefore conditional on this alignment and overstate what
would be expected under asynchronous drift in a real system. Miss ratio is also partly determined by the active policy, so the
CUSUM monitor can mistake a policy-induced change for workload change; the KS confirmation was designed
to address this, but in the frozen run it had only a small miss-ratio effect
(0.456 without KS vs.\ 0.453 with KS) despite being active on all 643 alarms.

\emph{Internal validity.} The controller configuration was selected by a sweep
on training traces before the fresh test set was touched. We mitigated the
resulting selection freedom by freezing all parameters and evaluating the test
traces once; nevertheless, the configuration was chosen to minimize training
miss ratio, and a reader should treat it as one degree of freedom.
Re-implementations of ARC, S3-FIFO, 2Q, LeCaR and the LPR-style
heuristic may differ from the original implementations, and baselines were not
tuned per workload, which can favour or penalize them. The simulator was
sanity-checked against known behaviour but not against reference
implementations. Seeds were fixed and the shadow simulator's sampling is
deterministic.

\emph{External validity.} A single-threaded trace-driven simulator omits
concurrency, memory pressure from multiple processes, reclaim interactions,
writeback, I/O queueing and NUMA effects. Cost figures come from an
unoptimized simulator on a shared virtual machine and should be read as
relative magnitudes. Results at one nominal capacity (with sensitivity at one other) and one trace
length may not transfer to other scales. Findings are
limited to the evaluated workload families.

\emph{Statistical conclusion validity.} The switching rule applies a normal
quantile ($z = 1.28$) to only four autocorrelated shadow blocks; a
$t$-quantile with three degrees of freedom is about 1.64, so the nominal 90\%
bound is optimistic. A post hoc seed-5001 robustness run with
$t_{0.90,3}$ gave DA miss 0.4586 [0.4282,0.4888] vs default 0.4529
[0.4224,0.4834]. KS $p$-values on dependent, tied reuse times are not calibrated error rates,
so the confirmation stage provides a heuristic gate rather than formal error
control; its ablation showed a small miss-ratio effect (0.456 vs.\ 0.453), and
the effective mechanism is better described as a CUSUM trigger with sampled
shadow evaluation. KS was enabled in the frozen run, the margin was $K = 0$, and no
validated two-stage statistical gate is claimed. Traces within a scenario share a generator,
so the effective number of independent workloads is smaller than 80, and the
stated $p$-values and intervals describe variability over generated traces,
not over real workloads. Wilcoxon tests with ten corrected comparisons have
limited power for differences of about 0.01 in mean miss ratio. No cost budget or agreement threshold was specified before the study. Several secondary results (for
for example the small KS ablation difference) are small or non-significant
differences and should not be read as evidence of no effect.

\section{Limitations and Scope of Claims}
\label{sec:limits}

This paper evaluates one controller design in a single, released simulator,
and it does not establish that the controller outperforms the best existing
replacement algorithms. In the reproducible main study the controller (0.453)
was 20.5\% below LRU (0.570) and better than S3-FIFO (0.468) on the mean, but
the S3-FIFO difference was not significant (adjusted $p = 0.09$), and
continuous shadow evaluation was significantly more accurate (0.423). The
controller cost 2.7 times LRU and 2.4 times S3-FIFO. The 76\% reduction in
shadow evaluations therefore came with a 0.030 loss in miss ratio in the main
study, not at equal miss ratio; only the real-program appendix shows a tie. We
make no claim of competitive cost or quality.

The present study is reproducible from the released package, but it remains a
trace-driven simulation: no in-kernel page-fault or reclaim latency was measured. CAR,
LIRS, CLOCK-Pro, CACHEUS, a 3L-style policy, an MGLRU-style policy and
set-dueling DIP are present in the released simulator and CSVs but are not
reported in the main tables; S4-FIFO, EACLOCK and Clock2Q+ were not implemented.
We do not claim the MGLRU-style implementation as a faithful baseline: a
page-level trace replay has no kernel aging or reclaim-pressure signal, so it
reduces to generation-indexed LRU rather than Linux MGLRU [38]. The frozen
controller used the five-candidate set including S3-FIFO; adding S3-FIFO
improved the miss ratio but did not make DA clearly better than the static
policy.

The KS confirmation had a small miss-ratio effect in the main study (0.456
without KS vs.\ 0.453 with KS) and was active on all 643 alarms. The cost
margin and dwell rule were inactive in the frozen configuration ($K = 0$, zero
dwell blocks), so their ablations do not estimate their benefit. The recency
filter acted at 47 of 213 switches without changing the mean miss ratio. The
effective mechanism was a CUSUM trigger with sampled shadow evaluation, not a
validated two-stage statistical gate. The evidence
is generated, with phase boundaries aligned to epochs and loop-heavy workloads
overrepresented; both choices favour detection and MRU-like heuristics. The
real-program traces come from fifteen instrumented kernels, six of which have
miss ratios below 0.02 at $C = 1000$; fifteen program-level clusters give
limited power and the mean is dominated by five high-miss kernels where DA and
S3-FIFO are nearly tied. A faithful kernel MGLRU baseline is not available:
the reported MGLRU-style policy is a page-level abstraction. No in-kernel
page-fault or reclaim latency was measured; the cost model uses a user-space
minor-fault calibration. The public traces are storage block/cache traces, not
paging traces, and only a prefix of each was replayed. The four-cluster
bootstrap of the synthetic main study has only 35 distinct resamples and is
reported as a coarse sensitivity check only. The literature review is
narrative, not systematic, and its limitations in Table~I are partly inferred.

Post hoc measurements (Section~\ref{sec:posthoc}) and recorded post hoc follow-ups
(Section~\ref{sec:posthoc_followups}) characterize trigger variants, rankability,
detector alternatives, stronger baselines, and cost uncertainty. The oracle
trigger uses synthetic ground-truth boundaries and is not deployable; the FIU
oracle produced no confirmations because FIU trace files expose no
ground-truth bounds. PH/ADWIN/SIEVE/W-TinyLFU/LHD are independent simplified
implementations in the released simulator with documented assumptions; they
are not authors' reference implementations. Available real-trace follow-ups
used 6 FIU public split traces (500k-request prepared splits), 15-program
traces, and three original FIU traces downloaded from [39] at deterministic
1/5/10\% WSS capacities.

The main remaining gaps are external validity and system implementation.
Bounded 1M-reference prefixes of the other six PDL families (MSR Cambridge,
Alibaba, Tencent, CloudPhysics, Twitter and MetaKV) were downloaded and run in
Phase 6; the [39] project page, FTP listing, exact FIU trace filename, README
and site-wide non-commercial copyright notice were verified, but no
dataset-specific licence file was found, so use is reported as non-commercial
research with citation rather than as use under a dataset licence. A
kernel-faithful MGLRU baseline and an in-kernel replacement prototype were not
built. A real userspace asynchronous shadow-evaluation prototype was built and
measured: across 1{,}232 completed evaluations the worker spent 577.8\,$\mu$s
wall and 575.5\,$\mu$s CPU per evaluation, and at the frozen model's ratio
$R=783$ this implies a 0.74\,$\mu$s miss-cost break-even, below the
1.44--4.06\,$\mu$s interval and the 1.98--2.06\,$\mu$s minor-fault calibration.
It is an unoptimised userspace measurement, not kernel or optimised latency,
and the prototype was worse than synchronous DA. In-kernel fault-path
durations were measured in a QEMU guest, but reclaim latency was not reached.
The cost calibration is a user-space minor-fault microbenchmark. These
limitations bound the claims to the evaluated simulator, traces, and calibration.

\section{Conclusion}
\label{sec:conclusion}

This paper asks when drift gating pays off in page replacement and answers it
as a reproducible measurement study. The evidence includes the frozen main
study, post hoc diagnostics, and post hoc follow-ups: a CUSUM-triggered
controller with sampled shadow evaluation is compared against continuous
shadow evaluation, static baselines, trigger alternatives, PH/ADWIN detectors,
and SIEVE/W-TinyLFU baselines; a cost break-even model is calibrated by a
user-space minor-fault measurement; and code, seeds, outputs, and protocols are
released.

On 80 non-stable synthetic traces, gating saved 76\% of shadow evaluations
but had higher miss ratio than continuous evaluation (0.453 vs.\ 0.423) and
was not significantly different from S3-FIFO (Holm-adjusted $p=0.09$). The FIU
capacity-sampling fault remained a failed external validation after repair.
Post hoc follow-ups showed a sharp rankability boundary without a
meaningful FIU repair gain, CUSUM budget-frontier advantage over
periodic/random but not over oracle/ContShadow accuracy, and no oracle-gap
closure by selected PH/ADWIN. On available real traces with SIEVE/W-TinyLFU,
DA showed no accuracy advantage. The 2.2\,$\mu$s minor-fault calibration and
illustrative miss-cost/async variants locate the DA--ContShadow crossover, but
the total-time envelope shows DA is never best in these measurements: CLOCK,
S3-FIFO and ContShadow each occupy part of the lower envelope and DA remains
at least 3.27\,ns/reference above the best alternative. External FIU originals
from [39] reproduced the absence of a DA advantage; stronger claims require the
remaining production traces and kernel-level measurements. A separate
hindsight-oracle study on 78 real/public traces found no replicated headroom
over the best static policy: 0/117 charged family-capacity groups at
50{,}000-reference windows, 1/117 at 100{,}000, and a ten-record extension
retained \texttt{PATH\_A\_STANDS}.

\section*{Appendix: Exploratory Real-Program Evidence}
\addcontentsline{toc}{section}{Appendix: Exploratory Real-Program Evidence}

\emph{Uncertainty note.} The trace-level interval is the primary uncertainty
statement throughout. As a sensitivity note only, the four-cluster bootstrap
on the synthetic traces has $\binom{7}{4} = 35$ distinct cluster resamples and
gives DA--S3-FIFO $-0.0154$ [$-0.0294$, $-0.0048$], excluding zero where the
primary trace-level paired interval does not. With only 35 resamples this is a
coarse check, not reliable evidence either way.

\emph{Real programs and whole-program holdout.} We recorded page traces by
instrumenting fifteen real programs: an index tree, a chained hash join, a CSR
breadth-first search, an in-place quicksort of 20-byte records, a column-major
Jacobi stencil, sparse-matrix value gathers, a binary heap with random
decrease-key, a histogram scatter over a 16 MB key array, a random walk over a
CSR adjacency array, an open-addressing hash table, a column-major
matrix-vector product over a 32 MiB matrix, pointer chasing through a 64 MiB
linked list, random-window string search over a 16 MiB text, union-find with
path compression over four million elements, and binary search over a 32 MiB
sorted key array. Four seeds were recorded per program (300,000 references
each; $C = 1000$, 100,000 warm-up). To remove seed-level leakage, we held out
\emph{whole programs}: for each of fifteen folds,
all tunable baselines and DA were tuned on the other fourteen programs and
evaluated once on the four traces of the omitted program. We also added
principle-level implementations of CAR, LIRS, CLOCK-Pro, CACHEUS, a 3L-style
policy, set-dueling DIP [37] and an MGLRU-style policy [38]; assumptions are
marked in the released source. The MGLRU-style policy is a page-level
abstraction of Linux MGLRU, not a kernel-faithful baseline: trace replay has no
kernel aging or reclaim-pressure signal.



The whole-program holdout means (Table~S.III in the supplementary material) are summarized here: The comparisons
below are exploratory; percentile intervals, Wilcoxon p-values, and Holm
adjustment are generated from the fifteen whole-program folds by the released
\texttt{appendix\_intervals.py} (bootstrap seed 7, 10,000 resamples) and are
fragile with fifteen clusters and many ties. With per-fold tuning and
whole-program holdout DA scored 0.3304, 6.7\% below LRU (0.3540). A post hoc
capacity-sampling repair ($f{=}1.0$, fixed a priori; \texttt{--cs-cap-f 1.0})
evaluated with the same whole-program holdout protocol produced identical mean
miss ratios to the original controller (DA 0.3304, S3-FIFO 0.3291, LRU 0.3540).



After Holm correction across these ten comparisons no difference remains
significant (smallest adjusted $p=0.096$ for 2Q). DA was not better than
S3-FIFO on average, so these fifteen-program results do not establish
superiority over the best static policy. Six of the fifteen kernels have miss ratios below 0.02 at this
capacity (hash join, breadth-first search, sort, stencil, radix scatter and
string search), and the mean is dominated by five high-miss kernels (sparse
gathers 0.72, random walk 0.83, hash table 0.71, linked list 0.94 and
union-find 0.73), where DA and S3-FIFO are nearly tied. In this independent
simulator, measured on the same 60 traces with the default frozen
configuration, DA cost 59.5 ns/reference against 27.5 for LRU and 24.7 for
S3-FIFO (2.2 and 2.4 times); the main-study 2.7-times figure differs because
the implementations and workloads differ, and on both DA is slower than
S3-FIFO. These are instrumented kernels, not production traces; the appendix
does not establish superiority or measure kernel latency. Source, seeds,
configurations, raw CSV output and analysis scripts, including
\texttt{appendix\_intervals.py}, are released in the accompanying reproduction
package.

\section*{Declarations}
\addcontentsline{toc}{section}{Declarations}

\emph{Data and code availability.} The simulator, trace generators, seeds,
configuration files, raw per-trace output and analysis scripts for the
reproducible main study are released in the accompanying
\texttt{reproduction/} package. \texttt{sweep\_selection.py} regenerates the
parameter-selection sweep, \texttt{run\_main\_study.sh} regenerates every
primary table and figure, \texttt{appendix\_intervals.py} regenerates the
appendix percentile intervals, and \texttt{cost\_model.py} together with
\texttt{prototype\_fault\_latency.cpp} regenerates the break-even analysis, and
\texttt{phase5/total\_time\_envelope.py} regenerates the total-time table,
figure and envelope.
The public-trace converter, the public-trace result and detection CSVs, the
independent capacity-pilot CSVs, and the boundary-offset detection CSVs are included in the
same package; the source traces are available from the repository in [39].
The instrumented programs and raw output for the exploratory real-program
appendix are included in the same package. No in-kernel prototype was built.

\emph{Funding.} This research was conducted independently without external
financial funding.

\emph{Conflict of interest.} M. Qasim is affiliated with both the University of
Management and Technology (UMT), Lahore, and Anoneurx Research Lab, Anoneurx;
the corresponding-author e-mail is at Anoneurx. The other authors are
affiliated with UMT. Anoneurx Astra, a tool developed at Anoneurx, was used for
prose drafting and code/analysis assistance, as disclosed below. Anoneurx
provided no funding for this work and had no role in study design, data
collection, analysis, interpretation or publication. No other competing
financial or non-financial interests are declared.

\emph{Use of AI-assisted tools.} AI-assisted tools, including Anoneurx Astra
and ChatGPT, were used to help draft and edit prose and to assist with code
development and analysis scripts. All experimental results were produced and
verified by the authors with the released code, and the authors take full
responsibility for the content of the paper.

\begin{thebibliography}{39}

\bibitem{r1} N. Megiddo and D. S. Modha, ``ARC: A self-tuning, low overhead
replacement cache,'' in \emph{Proc. 2nd USENIX Conf. File Storage Technol.
(FAST)}, 2003, pp. 115--130.

\bibitem{r2} S. Bansal and D. S. Modha, ``CAR: Clock with adaptive
replacement,'' in \emph{Proc. 3rd USENIX Conf. File Storage Technol. (FAST)},
2004, pp. 187--200.

\bibitem{r3} S. Jiang and X. Zhang, ``LIRS: An efficient low inter-reference
recency set replacement policy to improve buffer cache performance,'' in
\emph{Proc. ACM SIGMETRICS}, 2002, pp. 31--42.

\bibitem{r4} T. Johnson and D. Shasha, ``2Q: A low overhead high performance
buffer management replacement algorithm,'' in \emph{Proc. 20th Int. Conf. Very
Large Data Bases (VLDB)}, 1994, pp. 439--450.

\bibitem{r5} J. Yang, Y. Zhang, Z. Qiu, Y. Yue, and K. V. Rashmi, ``FIFO
queues are all you need for cache eviction,'' in \emph{Proc. 29th ACM Symp.
Oper. Syst. Princ. (SOSP)}, 2023, pp. 130--149.

\bibitem{r6} J. Gama, I. \v{Z}liobait\.{e}, A. Bifet, M. Pechenizkiy, and A.
Bouchachia, ``A survey on concept drift adaptation,'' \emph{ACM Comput.
Surv.}, vol. 46, no. 4, Art. no. 44, 2014.

\bibitem{r7} A. Bifet and R. Gavald\`a, ``Learning from time-changing data
with adaptive windowing,'' in \emph{Proc. SIAM Int. Conf. Data Mining (SDM)},
2007, pp. 443--448.

\bibitem{r8} G. Glass and P. Cao, ``Adaptive page replacement based on memory
reference behavior,'' in \emph{Proc. ACM SIGMETRICS}, 1997, pp. 115--126.

\bibitem{r9} Y. Smaragdakis, S. Kaplan, and P. Wilson, ``EELRU: Simple and
effective adaptive page replacement,'' in \emph{Proc. ACM SIGMETRICS}, 1999,
pp. 122--133.

\bibitem{r10} I. Ari, A. Amer, R. Gramacy, E. L. Miller, S. A. Brandt, and D.
D. E. Long, ``ACME: Adaptive caching using multiple experts,'' in \emph{Proc.
Workshop Distrib. Data Struct. (WDAS)}, 2002.

\bibitem{r11} K. Vietri \emph{et al.}, ``Driving cache replacement with
ML-based LeCaR,'' in \emph{Proc. 10th USENIX Workshop Hot Topics Storage File
Syst. (HotStorage)}, 2018.

\bibitem{r12} L. V. Rodriguez \emph{et al.}, ``Learning cache replacement
with CACHEUS,'' in \emph{Proc. 19th USENIX Conf. File Storage Technol. (FAST)},
2021, pp. 341--354.

\bibitem{r13} C. Waldspurger, T. Saemundsson, I. Ahmad, and N. Park, ``Cache
modeling and optimization using miniature simulations,'' in \emph{Proc. USENIX
Annu. Tech. Conf. (ATC)}, 2017, pp. 487--498.

\bibitem{r14} C. A. Waldspurger, N. Park, A. Garthwaite, and I. Ahmad,
``Efficient MRC construction with SHARDS,'' in \emph{Proc. 13th USENIX Conf.
File Storage Technol. (FAST)}, 2015, pp. 95--110.

\bibitem{r15} W. Zhou, Z. Niu, Y. Xiong, J. Fang, and Q. Wang, ``3L-Cache:
Low overhead and precise learning-based eviction policy for caches,'' in
\emph{Proc. 23rd USENIX Conf. File Storage Technol. (FAST)}, 2025, pp.
237--254.

\bibitem{r16} H. Kim, ``Learning-based page replacement scheme for efficient
I/O processing,'' \emph{Sci. Rep.}, vol. 15, 2025, doi:
10.1038/s41598-025-88736-4.

\bibitem{r17} G. Yuan, X. Li, S. Wu, Y. Zhu, D. Yang, C. Yao, X. Tang, and G.
Chen, ``EACLOCK: An efficient and adaptive page replacement algorithm to
understand your workload,'' \emph{IEEE Trans. Knowl. Data Eng.}, vol. 38,
no. 11, pp. 7387--7401, Nov. 2026, doi: 10.1109/TKDE.2026.3724159.
\url{https://doi.org/10.1109/TKDE.2026.3724159}.

\bibitem{r18} Y. Zhai \emph{et al.}, ``Clock2Q+: A simple and efficient
replacement algorithm for metadata cache in VMware vSAN,'' \emph{Proc. VLDB
Endow.}, vol. 19, no. 12, pp. 3874--3887, 2026, doi: 10.14778/3827998.3828002.

\bibitem{r19} H. Xia, W. Nixon, B. D. Marthen, P. Bhandari, and J. Yang,
``Learning-augmented heuristics: Simple yet smart, robust and interpretable
cache eviction,'' in \emph{Proc. 20th USENIX Symp. Oper. Syst. Design Implement.
(OSDI)}, 2026, pp. 2241--2259.

\bibitem{r20} P. Lu, J. Lu, A. Liu, and G. Zhang, ``Early concept drift
detection via prediction uncertainty,'' in \emph{Proc. AAAI Conf. Artif.
Intell.}, 2025; arXiv:2412.11158.

\bibitem{r21} N. Tan, Y.-C. Shih, D. Yang, and A. Salunkhe, ``Flexible and
efficient drift detection without labels,'' arXiv:2506.08734, 2025.

\bibitem{r22} P. Lu, J. Lu, A. Liu, E. Yu, and G. Zhang, ``Autonomous concept
drift threshold determination,'' arXiv:2511.09953, 2025.

\bibitem{r23} N. Harshit and K. Mounvik, ``Improving real-time concept drift
detection using a hybrid transformer-autoencoder framework,''
arXiv:2508.07085, 2025.

\bibitem{r24} S. Bandyopadhyay, J. Bose, and S. Roy Chowdhury, ``A hybrid
framework for real-time data drift and anomaly identification using
hierarchical temporal memory and statistical tests,'' \emph{Int. J. Math. Eng.
Manag. Sci.}, vol. 10, no. 3, pp. 777--796, 2025, doi:
10.33889/ijmems.2025.10.3.039.

\bibitem{r25} D. N. Assis and V. M. A. Souza, ``ADWIN-U: Adaptive windowing
for unsupervised drift detection on data streams,'' \emph{Knowl. Inf. Syst.},
2025, doi: 10.1007/s10115-025-02523-1.

\bibitem{r26} W. Shen, Y. Chen, R. Chen, and H. Chen, ``Towards fully-fledged
GPU multitasking via proactive memory scheduling,'' arXiv:2512.24637, 2025.

\bibitem{r27} Y. Zheng \emph{et al.}, ``gpu\_ext: Extensible OS policies for
GPUs via eBPF,'' arXiv:2512.12615, 2025.

\bibitem{r28} Y. Liu, M. Yu, and J. Yang, ``When fancy eviction fails:
Rethinking cache replacement for LLM prefix reuse,'' arXiv:2609.28870, 2026,
doi: 10.48550/arXiv.2609.28870.

\bibitem{r29} N. Patel and H. Maniar, ``Comparative evaluation of cache
replacement algorithms for time-bounded stale-tolerant data workloads,''
\emph{Int. J. Sci. Res. Technol.}, vol. 3, no. 6, pp. 556--564, 2026, doi:
10.5281/zenodo.20576370. \url{https://doi.org/10.5281/zenodo.20576370};
\href{https://ijsrtjournal.com/article/comparative-evaluation-of-cache-replacement-algorithms-for-time-bounded-stale-tolerant-data-workloads}{publisher landing page}.

\bibitem{r30} E. S. Page, ``Continuous inspection schemes,'' \emph{Biometrika},
vol. 41, no. 1--2, pp. 100--115, 1954.

\bibitem{r31} D. Narayanan, A. Donnelly, and A. Rowstron, ``Write
off-loading: Practical power management for enterprise storage,'' \emph{ACM
Trans. Storage}, vol. 4, no. 3, Art. no. 10, 2008.

\bibitem{r32} L. A. Belady, ``A study of replacement algorithms for a
virtual-storage computer,'' \emph{IBM Syst. J.}, vol. 5, no. 2, pp. 78--101,
1966.

\bibitem{r33} F. Wilcoxon, ``Individual comparisons by ranking methods,''
\emph{Biometrics Bull.}, vol. 1, no. 6, pp. 80--83, 1945.

\bibitem{r34} T. Azim, A. Nadeem, and A. Ailamaki, ``Adaptive cache mode
selection for queries over raw data,'' in \emph{Proc. 9th Int. Workshop
Accelerating Analytics Data Manag. Syst. Using Modern Processor Storage Archit.
(ADMS@VLDB)}, 2018, pp. 51--65.

\bibitem{r35} S. Holm, ``A simple sequentially rejective multiple test
procedure,'' \emph{Scand. J. Stat.}, vol. 6, no. 2, pp. 65--70, 1979.

\bibitem{r36} S. Jiang, F. Chen, and X. Zhang, ``CLOCK-Pro: An effective
improvement of the CLOCK replacement,'' in \emph{Proc. USENIX Annu. Tech. Conf.
(ATC)}, 2005, pp. 323--336.

\bibitem{r37} M. K. Qureshi, A. Jaleel, Y. N. Patt, S. C. Steely Jr., and J.
Emer, ``Adaptive insertion policies for high performance caching,'' in
\emph{Proc. 34th Int. Symp. Comput. Archit. (ISCA)}, 2007, pp. 381--391, doi:
10.1145/1250662.1250709.

\bibitem{r38} Linux Kernel Organization, ``Multi-gen LRU,'' Linux kernel
documentation, 2026. [Online]. Available:
\url{https://docs.kernel.org/mm/multigen_lru.html}

\bibitem{r39} CMU Parallel Data Laboratory, ``Cache eviction tactics:
cacheDatasets trace collection,'' online data repository and format README,
accessed Nov. 24, 2026. The project page links to the \texttt{cacheDatasets}
repository; its FIU directory contains
\texttt{fiu\_madmax-110108-112108.oracleGeneral.zst}. The README specifies
zstd-compressed 24-byte request records and citation requirements, but no
dataset-specific licence; PDL's site-wide copyright notice says materials
distributed by the server are provided on a non-commercial basis. [Online].
Available: \url{https://pdl.cmu.edu/CacheEviction/index.shtml};
\url{https://ftp.pdl.cmu.edu/pub/datasets/twemcacheWorkload/cacheDatasets/};
\url{https://ftp.pdl.cmu.edu/pub/datasets/twemcacheWorkload/cacheDatasets/README};
\url{https://pdl.cmu.edu/Publications/copyright.shtml}.

\end{thebibliography}

\end{document}
